\documentclass[11pt]{article}
\usepackage[letterpaper,margin=1in]{geometry}
\usepackage[T1]{fontenc}
\usepackage[utf8]{inputenc}
\usepackage{lmodern,microtype}
\usepackage{amsmath,amssymb,mathtools,bm,slashed,mathrsfs}
\usepackage{xcolor,enumitem,booktabs,array,needspace}
\usepackage[most]{tcolorbox}
\usepackage{fancyhdr,titlesec,parskip,hyperref}
\definecolor{qftblue}{HTML}{245A8D}
\definecolor{qftgreen}{HTML}{2F6B4F}
\definecolor{qftgray}{HTML}{F4F6F8}
\hypersetup{colorlinks=true,linkcolor=qftblue,urlcolor=qftblue,pdfauthor={Nima Lashkari}}
\pagestyle{fancy}\fancyhf{}
\lhead{PHYS 662: Quantum Field Theory I}\cfoot{\thepage}
\setlength{\headheight}{14pt}
\titleformat{\section}{\Large\bfseries\color{qftblue}}{\thesection.}{0.6em}{}
\titleformat{\subsection}{\large\bfseries\color{qftgreen}}{\thesubsection}{0.6em}{}
\setlist[enumerate]{leftmargin=2em,itemsep=0.25em,topsep=0.35em}
\setlist[itemize]{leftmargin=2em,itemsep=0.2em,topsep=0.35em}
\newcommand{\dd}{\mathrm d}
\newcommand{\ii}{\mathrm i}
\newcommand{\eps}{\varepsilon}
\newcommand{\tr}{\operatorname{tr}}
\newcommand{\Tr}{\operatorname{Tr}}
\newcommand{\diag}{\operatorname{diag}}
\newcommand{\Lag}{\mathcal L}
\newcommand{\D}{\mathcal D}
\newcommand{\vac}{\lvert0\rangle}
\newcommand{\vbra}{\langle0\rvert}
\newcommand{\Tord}{\mathrm T}
\newcommand{\ket}[1]{\lvert#1\rangle}
\newcommand{\bra}[1]{\langle#1\rvert}
\newtcolorbox{keyidea}[1][]{enhanced,breakable,colback=qftgray,colframe=qftblue,
 boxrule=0.7pt,arc=1.5mm,title={#1},fonttitle=\bfseries,coltitle=white}
\newtcolorbox{remarkbox}[1][]{enhanced,breakable,colback=white,colframe=qftgreen,
 boxrule=0.6pt,arc=1.5mm,title={#1},fonttitle=\bfseries,coltitle=white}
\newcommand{\intp}{\int\!\frac{\dd^4p}{(2\pi)^4}}
\newcommand{\MSbar}{\overline{\mathrm{MS}}}
\newcommand{\module}[1]{\clearpage\section{#1}}
\newcommand{\braket}[2]{\langle#1\mid#2\rangle}
\newcommand{\dPi}{\dd\Pi}
\newcommand{\cM}{\mathcal M}
\newcommand{\cH}{\mathcal H}
\newcommand{\cS}{\mathcal S}
\newcommand{\Zp}{Z_{\rm pole}}
\newcommand{\Kallen}{\mathcal K}
\rhead{Homework Set 8}
\hypersetup{pdftitle={PHYS 662 Homework Set 8: Scattering, Rates, and LSZ}}
\newlist{parts}{enumerate}{1}
\setlist[parts]{label=\textbf{(\alph*)},leftmargin=2.3em,itemsep=0.65em,topsep=0.5em}
\newcommand{\problem}[2]{\clearpage\section*{\color{qftblue}Problem #1: #2}}
\newcommand{\connection}[1]{{\small\textit{Lecture connection: #1}\par}\medskip}
\begin{document}
\begin{center}
{\scshape\Large PHYS 662 --- Quantum Field Theory I\par}
\vspace{0.3em}{\Huge\bfseries\color{qftblue}Homework Set 8\par}
\vspace{0.3em}{\large Scattering States, Rates, and LSZ Reduction\par}
\vspace{0.3em}{\large Lectures 23--25 \quad Fall 2026\par}
\end{center}\medskip
\textbf{Name:} \rule{0.43\textwidth}{0.4pt}\hfill
\textbf{Due:} \rule{0.24\textwidth}{0.4pt}

\textbf{Instructions.} Complete all four problems, worth 25 points each.
Show the origin of normalization and identical-particle factors. Keep
physical masses distinct from arbitrary Lagrangian parameters. Explain
the assumptions needed for every asymptotic limit. You may use results
from earlier parts, including when you have not completed their derivation.
No loop integral evaluation is required.

\begin{keyidea}[Common conventions]
Use \(\hbar=c=1\), \(\eta=\diag(1,-1,-1,-1)\), and
\[
 \langle p,s|q,r\rangle=2E_p(2\pi)^3\delta^{(3)}(\bm p-\bm q)\delta_{sr},
 \qquad \dPi_p=\frac{\dd^3p}{(2\pi)^3\,2E_p}.
\]
The connected scattering amplitude is defined by
\[
 \langle f|S-1|i\rangle_c
 =\ii(2\pi)^4\delta^{(4)}(P_f-P_i)\cM_{fi},\qquad S=1+\ii T.
\]
A free scalar propagator includes its numerator \(\ii\):
\[
 D_F(p)=\frac{\ii}{p^2-m^2+\ii0},\qquad
 (\Box+m^2)D_F(x)=-\ii\delta^{(4)}(x).
\]
For a stable particle, \(\langle\Omega|\phi(0)|p\rangle=\sqrt{\Zp}\)
and \(\widetilde G_2(p)\sim\ii\Zp/(p^2-M^2+\ii0)\).
Use \(\cS_f=\prod_a n_a!\) when integrating over all labeled final momenta,
and
\[
 \dd\Phi_n(P)=(2\pi)^4\delta^{(4)}\!\left(P-\sum_jp_j\right)\prod_j\dPi_{p_j}.
\]
The invariant incoming flux is
\(F=4\sqrt{(p_A\cdot p_B)^2-m_A^2m_B^2}\).
For two-body kinematics use
\(\Kallen(x,y,z)=x^2+y^2+z^2-2xy-2xz-2yz\).
\end{keyidea}

\textbf{Coverage.} Spectral weight and scattering states; two-body decay
and resonances; cross sections and unitarity; scalar and fermion reduction.
All masses are positive. External particles are stable except for the
narrow parent in Problem 2. Assume weak couplings and no additional lighter
bound states where needed.

\problem{1}{What is a particle, and on what space is the S-matrix unitary?}
\connection{Lecture 23, Sections 1--6 and 8.}
Consider a Hermitian scalar \(\phi\) with \(\langle\phi\rangle=0\) in a
positive-metric Hilbert space, with a stable isolated mass \(M\).
\begin{parts}
\item \textbf{(4 points)} Insert exact energy-momentum states into
\(\langle\Omega|\phi(x)\phi(0)|\Omega\rangle\). With the convention
\[
 \sum_n\int\dd\mu(n)\,(2\pi)^4\delta^{(4)}(q-p_n)
 |\langle\Omega|\phi(0)|n\rangle|^2=2\pi\theta(q^0)\rho(q^2),
\]
derive the spectral representation of the Feynman two-point function
and explain why \(\rho\geq0\). Why is it not just a count of states?

\item \textbf{(4 points)} Evaluate the one-particle contribution using
the stated covariant normalization. Relate its delta-function coefficient
to the propagator pole residue. Under \(\phi\to c\phi\), with \(c>0\),
which of the mass, residue, and spectral weight change?

\item \textbf{(4 points)} In unbroken massive \(\lambda\phi^4\) theory,
explain why the first multiparticle continuum in \(\rho_\phi\) starts at
\(9M^2\), whereas the even operator
\(\phi^2-\langle\phi^2\rangle\) can have a continuum at \(4M^2\).
State the needed spectral assumptions. Why is the four-point bubble's
cut at \(s=4M^2\) consistent with these answers?

\item \textbf{(4 points)} Derive the direct and exchange terms in the
inner product of two identical scalar states. Verify that
\(\tfrac1{2!}\int\dPi_p\dPi_q\,|p,q\rangle\langle p,q|\)
acts as the identity on the two-particle sector. Explain why a local
product \(\phi(x)\phi(y)\ket{\Omega}\) is not automatically in that sector.

\item \textbf{(5 points)} Let \(\Omega_{\rm in/out}\) be isometries from
\(\cH_{\rm as}\) onto a common scattering subspace. State the wave-packet
comparison condition, including the two time directions, and prove that
\(S=\Omega_{\rm out}^\dagger\Omega_{\rm in}\) is unitary. Distinguish
equality of the two scattering ranges from asymptotic completeness in
the full Hilbert space. Why is physical switching off of the coupling
not required?

\item \textbf{(4 points)} For an elastic projector \(P\), show that
\((PSP)^\dagger(PSP)=P-PS^\dagger(1-P)SP\). Interpret the missing probability.
Then derive \(-\ii(T-T^\dagger)=T^\dagger T\) and explain why a nonzero real
tree transition amplitude requires an imaginary part at a higher order.
\end{parts}

\problem{2}{A two-body decay and its resonance pole}
\connection{Lecture 24, Sections 2--4 and 7.}
Take \(\Lag_{\rm int}=-gA\chi^2/2\), with \(M_A>2m_\chi\), real small
\(g\), and stable \(\chi\) particles. Assume additional interactions
stabilize the potential without contributing to this leading decay.
Treat \(A\) as a narrow prepared excitation, not an exact asymptotic state.
\begin{parts}
\item \textbf{(4 points)} Explain how the regulated limit
\[
 \frac1{\mathcal T}\left|\int_{-\mathcal T/2}^{\mathcal T/2}
 \dd t\,e^{\ii\omega t}\right|^2\longrightarrow2\pi\delta(\omega)
\]
and the norms \(2EV\) lead to
\(\dd\Gamma=|\cM|^2\dd\Phi_2/(2M_A\cS_f)\).
Show where the volume and observation time disappear.

\item \textbf{(5 points)} Integrate the two-body momentum delta functions
in the rest frame of \(A\). Derive \(\dd\Phi_2/\dd\Omega\), the momentum
magnitude \(p_*\), and the integrated phase space. Identify the threshold
condition rather than inferring it from a square root alone.

\item \textbf{(4 points)} Derive the vertex rule, the tree amplitude,
and the total partial width \(\Gamma_{A\to\chi\chi}\). Explain why the
\(1/2\) in the Lagrangian does not remove the need for a final-state
\(1/2!\). Check mass dimensions and the threshold limit.

\item \textbf{(4 points)} If this is the only open channel, find the mean
proper lifetime and half-life. Give the laboratory mean lifetime for a
parent packet sharply peaked at energy \(E_A\). If another distinct channel
has width \(\Gamma_Y\), state the new total lifetime and branching fraction
to \(\chi\chi\).

\item \textbf{(4 points)} Starting from
\(\widetilde G_A(s)\simeq\ii\mathcal Z_R/(s-M_A^2+\ii M_A\Gamma_A)\),
locate the positive-energy pole to first order in \(\Gamma_A/M_A\).
Determine the damping of the rest-frame amplitude and of its modulus
squared. Explain the factor of two and why this is not an exact stable
K\"all\'en--Lehmann delta-function pole.

\item \textbf{(4 points)} Prove the narrow-width identity
\[
 \frac1{(s-M_A^2)^2+M_A^2\Gamma_A^2}
 \longrightarrow\frac{\pi}{M_A\Gamma_A}\delta(s-M_A^2)
\]
under integration against a smooth function. State two circumstances in
which using this identity for an observed line shape is unreliable.
\end{parts}

\problem{3}{Two stable species: flux, cross sections, and unitarity}
\connection{Lecture 24, Sections 4--6 and 8.}
Consider two real scalar species with masses \(m_\chi>m_\phi>0\) and
\(\Lag_{\rm int}=-g\phi^2\chi^2/4\), with \(g>0\) small. Independent
\(\mathbb Z_2\) parities make both species stable. Work at leading order,
and define \(\beta_a(s)=\sqrt{1-4m_a^2/s}\) in its physical region.
\begin{parts}
\item \textbf{(4 points)} Derive the four-leg vertex factor from Wick
contractions. Find the tree amplitudes for
\(\phi\phi\to\chi\chi\) and \(\phi\chi\to\phi\chi\).
Specify their physical thresholds in \(s\) and their final-state counting
factors. Why is there no extra incoming \(1/2!\) for two prepared beams?

\item \textbf{(4 points)} Starting from
\(F=4\sqrt{(p_A\cdot p_B)^2-m_A^2m_B^2}\), derive its center-of-mass and
stationary-target forms. Show that \(F/(4E_AE_B)\) equals
\(\sqrt{|\bm v_A-\bm v_B|^2-|\bm v_A\times\bm v_B|^2}\).
Explain why this flux factor need not be less than one.

\item \textbf{(5 points)} Calculate the differential and total event
cross sections for \(\phi\phi\to\chi\chi\) above threshold, displaying
the factor \(\beta_\chi/\beta_\phi\). Use the full labeled angular sphere.
Explain how the answer changes if you integrate over only a nonredundant
hemisphere without a symmetry factor.

\item \textbf{(4 points)} Calculate the total cross section for
\(\phi\chi\to\phi\chi\) at the same invariant \(s\) when both reactions
are allowed. Compare it to part (c), separating the effects of phase space
and identical-particle counting. In a stationary \(\chi\) target, express
\(s\) in terms of the incident \(\phi\) energy.

\item \textbf{(4 points)} Without evaluating a loop integral, use
unitarity to determine the \(\chi\chi\) intermediate-state contribution
to \(\operatorname{Im}\cM_{\phi\phi\to\phi\phi}(s,0)\) at order \(g^2\).
Check it independently against \(F\sigma_{\phi\phi\to\chi\chi}/2\).
Explain why the \(\chi\chi\) contribution vanishes below its threshold.

\item \textbf{(4 points)} A collider delivers integrated luminosity
\(\mathscr L_{\rm int}=\int\dd t\,\mathscr L(t)\) at fixed \(s\), with
uniform event efficiency \(\epsilon\). Express the expected detected
\(\chi\chi\) event count in terms of part (c). Contrast it with the
inclusive number of detected \(\chi\) particles at perfect acceptance.
Why is a cross section not a probability or a fraction of collisions?
\end{parts}

\problem{4}{Reduction from correlators and the choice of field}
\connection{Lecture 25, Sections 2--7.}
Use smooth wave packets before taking any plane-wave limit. In parts
(a)--(b) consider one stable Hermitian scalar of mass \(M\) and residue
\(\Zp\). In part (c) use the two-species interaction from Problem 3.
\begin{parts}
\item \textbf{(4 points)} For a positive-frequency Klein--Gordon solution
\(f\), define
\[
 a_f(t)=\frac{\ii}{\sqrt{\Zp}}\int\dd^3x\,
 f^*\overleftrightarrow{\partial_t}\phi.
\]
Derive its time derivative and the formula for
\(a_{\rm out}(f)-a_{\rm in}(f)\). State the meaning of the asymptotic
limits. Use Hermitian conjugation to show why incoming creation-operator
reduction also contributes \(+\ii/\sqrt{\Zp}\).

\item \textbf{(4 points)} Explain the origin of the spectator overlap
when one outgoing annihilation operator is replaced by an incoming one.
Derive the free identity \((\Box_x+M^2)D_F(x-y)=-\ii\delta^{(4)}(x-y)\)
from the equal-time commutator and derivatives of time ordering. Why would
using the homogeneous field equation alone give the wrong answer?

\item \textbf{(5 points)} Write the leading connected position-space
correlator with two \(\phi\) and two \(\chi\) insertions as a single vertex
integral. Apply scalar LSZ directly to obtain
\(\ii(2\pi)^4\delta^{(4)}(P_f-P_i)\cM_{\chi\chi,\phi\phi}\).
Track every \(\ii\), incoming/outgoing Fourier sign, and physical mass.
Then state the exact relation between this amplitude and the connected
kernel amputated with full propagators, allowing residues
\(Z_\phi\ne1\) and \(Z_\chi\ne1\).

\item \textbf{(4 points)} Demonstrate invariance of the scalar LSZ result
under \(\phi\to c\phi\), \(c>0\). For
\(O=(1+\alpha\Box)\phi\), determine its one-particle overlap and pole
residue. For which \(\alpha\) does this operator fail to interpolate the
particle? Explain why local contact terms do not repair that missing pole.

\item \textbf{(4 points)} For a stable Dirac field, list the external
spinors for an incoming fermion, outgoing fermion, incoming antifermion,
and outgoing antifermion. Explain where \(\sqrt{Z_\psi}\) enters if the
connected kernel has been amputated with full propagators. Distinguish
external spin sums from the average over an unpolarized incoming ensemble.

\item \textbf{(4 points)} For the Yukawa interaction \(-y\phi\bar\psi\psi\),
write the scalar-exchange tree amplitude for identical fermions:
\[
 \psi(p_1,s_1)\psi(p_2,s_2)\to\psi(p_3,s_3)\psi(p_4,s_4).
\] Include both contractions and their relative sign,
with unit residues at this order. Check antisymmetry under exchange of the
two final fermions. No spin trace or phase-space integral is required.
\end{parts}
\end{document}
